% Part:sets-functions-relations
% Chapter: size-of-sets
% Section: equinumerous-sets

\documentclass[../../../include/open-logic-section]{subfiles}

\begin{document}

\olfileid{sfr}{siz}{equ}

\olsection{等势}

我们对集合的“大小”有一个直观的概念，这对有限集来说是很适用的。
但对于无穷集呢？如果我们想要形式化地比较\emph{任意}大小的两个集合，一个很好的出发点是定义何时两个集合大小相同。以下是 Frege的论述：

\begin{quote}
  如果一名侍者想确认他在桌上摆放的餐刀恰好和餐盘一样多，
  他并不需要清点任何一个，
  只要他在每个盘子的右边放上一把刀，
  使得桌上的每把刀都位于某个盘子的右边即可。这样，
  餐盘与餐刀就通过这同样的空间关系而被一一对应起来了。
  \citep[\S70]{Frege1884}
\end{quote}

这段话的洞见可以通过以下形式定义来阐明：

\begin{defn}\ollabel{comparisondef}
  $A$ 与 $B$ 是\emph{等势的}，记作 $\cardeq{A}{B}$，
  当且仅当存在一个双射 $f \colon A \to B$。
\end{defn}

\begin{prop}\ollabel{equinumerosityisequi}
等势是一个等价关系。
\end{prop}

\begin{proof}
我们必须证明等势是自反的、对称的和传递的。设 $A$、$B$ 和 $C$ 为集合。

\emph{自反性。} 恒等映射 $\Id{A} \colon A \to A$，其中对所有 $x \in A$ 有 $\Id{A} (x) = x$，是一个双射。所以 $\cardeq{A}{A}$。

\emph{对称性。} 假设 $\cardeq{A}{B}$，即存在一个双射 $f\colon A \to B$。由于 $f$ 是双射，其逆 $f^{-1}$ 存在且也是双射。因此，$f^{-1}\colon B \to A$ 是一个双射，故 $\cardeq{B}{A}$。

\emph{传递性。} 假设 $\cardeq{A}{B}$ 且 $\cardeq{B}{C}$，即存在双射 $f\colon A \to B$ 和 $g\colon B \to C$。那么复合 $\comp{f}{g}\colon A \to C$ 是双射，因此 $\cardeq{A}{C}$。
\end{proof}

\begin{prop}
若 $\cardeq{A}{B}$，则 $A$ 是可数的当且仅当 $B$ 是可数的。
\end{prop}

\begin{editorial}
以下证明根据是否包含 \olref[enm]{sec} 节，分别使用 \olref[enm]{defn:enumerable} 或 \olref[enm-alt]{defn:enumerable}。
\end{editorial}

\begin{proof}
假设 $\cardeq{A}{B}$，即存在某个双射 $f \colon A \to B$，并假设 $A$ 是可数的。
\oliflabeldef{sfr:siz:enm:defn:enumerable}{
  那么要么 $A = \emptyset$，要么存在一个满射函数 $g\colon \PosInt \to A$。若 $A = \emptyset$，则也有 $B = \emptyset$（否则就会存在元素 $y \in B$ 但没有 $x \in A$ 使得 $g(x) = y$）。另一方面，若 $g\colon \PosInt \to A$ 是满射，则 $\comp{g}{f} \colon \PosInt \to B$ 是满射。为说明这一点，令 $y \in B$。由于 $f$ 是满射，存在 $x \in A$ 使得 $f(x) = y$。由于 $g$ 是满射，存在 $n \in \PosInt$ 使得 $g(n) = x$。于是，
\[
(\comp{g}{f})(n) = f(g(n)) = f(x) = y
\]
因此 $\comp{g}{f}$ 是满射。我们得到 $\comp{g}{f}$ 是 $B$ 的一个枚举，故 $B$ 是可数的。}
{
  那么要么 $A = \emptyset$，要么存在一个双射 $g$，其值域是 $A$ 而定义域是 $\Nat$ 或自然数的一个初始段。若 $A = \emptyset$，则也有 $B = \emptyset$（否则就会存在某个 $y \in B$ 而没有 $x \in A$ 使得 $g(x) = y$）。所以假设我们有双射 $g$。那么 $\comp{g}{f}$ 是一个双射，其值域为 $B$ 而定义域与 $g$ 的相同（即 $\Nat$ 或其初始段），因此 $B$ 是可数的。}

如果 $B$ 是可数的，我们只需在上述论证中以双射 $f^{-1}\colon B \to A$ 替代 $f$，便可得出 $A$ 也是可数的。
\end{proof}

\begin{prob}
证明：如果 $\cardeq{A}{C}$ 且 $\cardeq{B}{D}$，并且 $A \cap B = C \cap D = \emptyset$，那么 $\cardeq{A \cup B}{C \cup D}$。
\end{prob}

\begin{prob}
证明：如果 $A$ 是无穷且可数的，那么 $\cardeq{A}{\Nat}$。
\end{prob}

\end{document}